\documentclass[12pt,a4paper]{article}

\usepackage[utf8]{inputenc}
\usepackage[T1]{fontenc}
\usepackage[brazil]{babel}
\usepackage{geometry}
\usepackage{graphicx}
\usepackage{xcolor}
\usepackage{booktabs}
\usepackage{float}
\usepackage{hyperref}
\usepackage{enumitem}
\usepackage{array}
\usepackage{tabularx}

\geometry{margin=2.4cm}
\hypersetup{
    colorlinks=true,
    linkcolor=blue!50!black,
    urlcolor=blue!50!black,
    citecolor=blue!50!black
}
\setlist[itemize]{noitemsep,topsep=3pt}
\setlist[enumerate]{topsep=3pt}
\urlstyle{same}
\sloppy

\newcommand{\resultimage}[2]{%
    \IfFileExists{data/results/#1}{%
        \includegraphics[#2]{data/results/#1}%
    }{%
        \IfFileExists{../data/results/#1}{%
            \includegraphics[#2]{../data/results/#1}%
        }{%
            \fbox{\parbox[c][4.8cm][c]{0.86\textwidth}{\centering
            Arquivo \texttt{#1} não encontrado. Execute \texttt{make run}.}}%
        }%
    }%
}

\title{
    \textbf{Detecção e Reconstrução de Malha Viária\\
    a partir de Imagens de Satélite}\\
    \large Trabalho Prático -- Disciplina de Reconhecimento de Padrões
}

\author{
    Beatriz Vocurca Frade\\
    Johnatan Augusto Moreira do Carmo\\
    Marina Alves Resende
}

\date{\today}

\begin{document}

\maketitle

\begin{abstract}
Este relatório apresenta um sistema de visão computacional clássica para detectar vias em imagens de satélite, segmentar regiões candidatas a estradas e reconstruir uma representação topológica da malha viária. A solução combina pré-processamento, realce local de contraste, descritores de cor e textura, suporte de estruturas lineares, morfologia matemática, esqueletonização e extração de grafos. A abordagem foi escolhida por ser interpretável, reprodutível e executável sem treinamento supervisionado, o que permite discutir de forma transparente tanto os acertos quanto as limitações do método.
\end{abstract}

\tableofcontents
\newpage

\section{Introdução}

A detecção automática de vias em imagens de satélite é uma tarefa importante para atualização cartográfica, planejamento urbano, monitoramento de infraestrutura, logística, resposta a emergências e análise territorial. Em uma imagem aérea, estradas costumam aparecer como regiões alongadas, relativamente homogêneas e conectadas. Entretanto, essa aparência não é exclusiva de vias: telhados, estacionamentos, solo exposto, sombras, trilhas agrícolas e áreas pavimentadas podem produzir padrões visuais semelhantes.

O objetivo deste trabalho é desenvolver um sistema capaz de:

\begin{itemize}
    \item identificar pixels candidatos a vias;
    \item diferenciar visualmente vias asfaltadas e vias de terra;
    \item sobrepor a detecção na imagem original;
    \item reduzir a máscara segmentada a um esqueleto topológico;
    \item exportar um grafo com vértices e arestas da malha viária.
\end{itemize}

O foco não é prometer desempenho perfeito em qualquer imagem, mas apresentar uma solução tecnicamente fundamentada, executável, analisável e compatível com o escopo da disciplina. Como não há, neste projeto, uma base anotada com máscaras de referência, a avaliação é qualitativa e baseada em métricas auxiliares produzidas pelo próprio pipeline.

\section{Fundamentação Teórica}

\subsection{Reconhecimento de padrões em imagens}

Reconhecimento de padrões em imagens consiste em identificar regularidades visuais associadas a uma classe de interesse. No caso de malhas viárias, os padrões relevantes incluem continuidade espacial, formato alongado, baixa variação local de textura, contraste com vegetação, presença de cruzamentos e repetição de estruturas lineares.

Uma dificuldade central é que a classe ``via'' não é definida apenas por cor. Estradas asfaltadas podem ser cinzas, azuladas, escuras ou claras, enquanto estradas de terra podem variar entre tons amarelados, alaranjados e avermelhados. Por isso, a solução proposta combina múltiplos indícios visuais em vez de depender de um único limiar.

\subsection{Segmentação}

Segmentação é o processo de separar uma imagem em regiões semanticamente relevantes. Neste trabalho, a etapa principal é uma segmentação binária: cada pixel é classificado como candidato a via ou fundo. Em seguida, os pixels de via recebem uma classe visual: asfalto ou terra.

A segmentação usa critérios de cor, textura e estrutura. Essa combinação reduz erros comuns de métodos puramente baseados em cor, pois uma região marrom só é tratada como possível via de terra quando também apresenta suporte linear.

\subsection{Morfologia matemática}

Operações morfológicas processam a forma das regiões segmentadas. O fechamento conecta pequenas quebras e preenche falhas internas, enquanto a abertura remove ruídos isolados. Após essas operações, componentes conectados muito pequenos ou muito compactos são filtrados para diminuir falsos positivos.

\subsection{Esqueletonização e grafos}

A esqueletonização transforma uma máscara binária espessa em linhas de espessura unitária, preservando aproximadamente a conectividade original. Sobre esse esqueleto, o grau de cada pixel é estimado pela vizinhança 8-conectada:

\begin{itemize}
    \item grau 1: extremidade;
    \item grau 2: ponto intermediário de uma via;
    \item grau maior ou igual a 3: cruzamento ou bifurcação.
\end{itemize}

Agrupando pixels de extremidade e interseção, o sistema cria nós. Os caminhos entre esses nós formam as arestas do grafo. A saída final mantém o maior componente conectado do grafo extraído, aproximando o requisito de uma malha viária conectada.

\section{Metodologia}

A implementação foi organizada em duas camadas. O módulo \texttt{src/malha\_viaria.py} contém o pipeline executável e reutilizável. O notebook \texttt{notebook/deteccao\_malha\_viaria.ipynb} funciona como roteiro didático para executar o sistema, visualizar resultados e interpretar as saídas.

O fluxo completo pode ser executado pela linha de comando:

\begin{verbatim}
python3 -m src.malha_viaria \
  --input data/raw \
  --results data/results \
  --graphs data/graphs
\end{verbatim}

Quando \texttt{data/raw/} está vazia, o sistema gera uma imagem sintética de demonstração. Essa imagem não substitui validação em dados reais, mas garante que a entrega seja executável e que todas as etapas do pipeline produzam arquivos de saída.

\subsection{Pré-processamento}

Cada imagem é carregada em RGB e, se necessário, redimensionada para limitar o custo computacional. Em seguida, aplica-se:

\begin{itemize}
    \item filtro bilateral, para reduzir ruído preservando bordas;
    \item conversão para Lab, HSV e tons de cinza;
    \item CLAHE no canal de luminância, para melhorar contraste local;
    \item cálculo de desvio padrão local, usado como medida simples de textura.
\end{itemize}

O uso de diferentes espaços de cor facilita separar propriedades complementares: HSV destaca saturação e matiz; Lab permite realce de luminância; RGB preserva relações diretas entre canais.

\subsection{Segmentação de vias}

A máscara candidata é construída pela combinação de regras:

\begin{itemize}
    \item pixels acinzentados, pouco saturados e homogêneos são candidatos a asfalto;
    \item pixels de tons quentes, homogêneos e com predominância relativa de vermelho são candidatos a terra;
    \item regiões com excesso de verde são removidas como provável vegetação;
    \item sombras muito escuras são removidas;
    \item vias de terra exigem suporte de estruturas lineares para evitar confusão com solo exposto.
\end{itemize}

O suporte linear combina bordas Canny dilatadas em diferentes orientações com uma resposta simplificada baseada na Hessiana. A intenção é favorecer padrões compridos e estreitos, compatíveis com vias, em vez de manchas compactas.

\subsection{Classificação asfalto/terra}

A classificação entre asfalto e terra é heurística e tem finalidade principalmente visual. Pixels de via com baixa saturação ou pequena diferença entre canais RGB são rotulados como asfalto. Pixels com matiz quente e maior saturação são rotulados como terra. Essa etapa é útil para interpretar a segmentação, mas é sensível a iluminação, resolução e aparência do solo.

\subsection{Extração topológica}

Depois da segmentação, a máscara é esqueletonizada. O algoritmo calcula o grau de cada pixel do esqueleto, agrupa pixels de vértice em componentes conectados e percorre os caminhos entre componentes. Cada caminho encontrado é exportado como uma aresta com comprimento aproximado em pixels e geometria em coordenadas da imagem.

\section{Descrição do Algoritmo}

O pipeline implementado segue as etapas abaixo:

\begin{enumerate}[label=\textbf{Etapa \arabic*:}, leftmargin=*]
    \item listar imagens de entrada em \texttt{data/raw/};
    \item usar uma imagem sintética se a pasta de entrada estiver vazia;
    \item carregar e redimensionar a imagem;
    \item aplicar filtro bilateral e CLAHE;
    \item calcular descritores em RGB, HSV, Lab e tons de cinza;
    \item estimar suporte linear por bordas direcionais e resposta de segunda ordem;
    \item combinar regras para produzir a máscara candidata;
    \item aplicar fechamento, abertura e filtragem por componentes conectados;
    \item classificar vias em asfalto ou terra;
    \item gerar máscara, segmentação colorida e sobreposição;
    \item esqueletonizar a máscara;
    \item extrair nós e arestas do grafo;
    \item exportar imagens, JSON e GraphML.
\end{enumerate}

\section{Resultados}

\subsection{Saídas geradas}

Para cada imagem processada, o sistema produz os arquivos da Tabela~\ref{tab:saidas}.

\begin{table}[H]
    \centering
    \caption{Arquivos gerados pelo pipeline.}
    \label{tab:saidas}
    \begin{tabularx}{\textwidth}{>{\ttfamily}lX}
        \toprule
        Arquivo & Conteúdo \\
        \midrule
        *\_mascara.png & máscara binária de vias detectadas \\
        *\_segmentada.png & visualização por classe: asfalto em azul e terra em laranja \\
        *\_sobreposicao.png & vias detectadas sobre a imagem original \\
        *\_esqueleto.png & esqueleto usado para extração topológica \\
        *\_grafo.png & grafo desenhado sobre a sobreposição \\
        *\_grafo.json & nós, arestas, métricas e geometria dos segmentos \\
        *\_grafo.graphml & grafo em formato compatível com ferramentas externas \\
        \bottomrule
    \end{tabularx}
\end{table}

\subsection{Visualização do pipeline}

\begin{figure}[H]
    \centering
    \resultimage{demo_sintetico_segmentada.png}{width=0.78\textwidth}
    \caption{Segmentação colorida: asfalto em azul e terra em laranja.}
    \label{fig:segmentada}
\end{figure}

\begin{figure}[H]
    \centering
    \resultimage{demo_sintetico_sobreposicao.png}{width=0.78\textwidth}
    \caption{Sobreposição das vias detectadas na imagem de entrada.}
    \label{fig:sobreposicao}
\end{figure}

\begin{figure}[H]
    \centering
    \resultimage{demo_sintetico_esqueleto.png}{width=0.78\textwidth}
    \caption{Esqueleto extraído da máscara binária.}
    \label{fig:esqueleto}
\end{figure}

\begin{figure}[H]
    \centering
    \resultimage{demo_sintetico_grafo.png}{width=0.78\textwidth}
    \caption{Grafo extraído a partir do esqueleto da malha detectada.}
    \label{fig:grafo}
\end{figure}

\subsection{Métricas auxiliares}

Na execução sintética presente em \texttt{data/results/}, o pipeline registrou as métricas da Tabela~\ref{tab:metricas}. Elas não substituem métricas supervisionadas, mas ajudam a verificar coerência, volume de detecção e complexidade topológica.

\begin{table}[H]
    \centering
    \caption{Métricas auxiliares da execução de demonstração.}
    \label{tab:metricas}
    \begin{tabular}{lr}
        \toprule
        Métrica & Valor \\
        \midrule
        Razão de pixels classificados como via & 7,47\% \\
        Pixels de via & 44.037 \\
        Componentes na máscara & 31 \\
        Nós no maior componente do grafo & 18 \\
        Arestas no maior componente do grafo & 19 \\
        Componentes antes da seleção do maior grafo & 33 \\
        Pixels classificados como asfalto & 24.469 \\
        Pixels classificados como terra & 19.568 \\
        \bottomrule
    \end{tabular}
\end{table}

\section{Análise de Desempenho}

Como o projeto não inclui máscaras anotadas, a avaliação é qualitativa. Os critérios observados foram:

\begin{itemize}
    \item continuidade das vias segmentadas;
    \item preservação de cruzamentos e bifurcações;
    \item coerência visual da classificação asfalto/terra;
    \item redução de regiões compactas que não representam vias;
    \item qualidade do esqueleto como aproximação topológica;
    \item conectividade do grafo exportado.
\end{itemize}

O resultado demonstra que a combinação de cor, textura e estrutura linear é capaz de recuperar trechos principais da malha e gerar um grafo inspecionável. Ao mesmo tempo, a figura segmentada evidencia falsos positivos, especialmente em objetos com cor ou textura parecida com vias. Essa limitação é esperada em uma abordagem sem treinamento supervisionado e reforça a importância de análise crítica em vez de uma conclusão baseada apenas em imagens visualmente agradáveis.

\section{Dificuldades Encontradas}

O problema é difícil porque a aparência visual de uma via muda muito entre regiões, escalas, sensores e condições de iluminação. Estradas asfaltadas podem se confundir com telhados, estacionamentos e sombras. Estradas de terra podem se confundir com solo exposto, talhões agrícolas e terrenos sem vegetação.

Outro desafio é a extração topológica. Pequenas falhas na máscara podem quebrar uma via em vários componentes, enquanto ruídos conectados podem criar cruzamentos inexistentes. Por isso, o grafo resultante deve ser interpretado como uma aproximação da estrutura viária, não como uma reconstrução cartográfica perfeita.

\section{Possíveis Melhorias Futuras}

Melhorias relevantes incluem:

\begin{itemize}
    \item usar uma base anotada para medir precisão, revocação, F1-score e IoU;
    \item treinar uma rede de segmentação, como U-Net, DeepLab ou SegFormer;
    \item calibrar limiares por validação cruzada em múltiplas cidades;
    \item aplicar fechamento topológico para reconectar quebras pequenas;
    \item remover falsos positivos com descritores de forma ou classificador supervisionado;
    \item comparar a malha extraída com dados do OpenStreetMap;
    \item incorporar imagens multiespectrais, quando disponíveis.
\end{itemize}

\section{Conclusão}

O sistema desenvolvido mostra que técnicas clássicas de visão computacional podem formar uma solução completa para detectar vias, gerar visualizações intermediárias e reconstruir uma malha viária aproximada como grafo. A principal vantagem da abordagem é a interpretabilidade: cada etapa pode ser visualizada, explicada e ajustada.

A solução também deixa claras suas limitações. Sem dados anotados e sem aprendizado supervisionado, o método depende de hipóteses visuais que podem falhar em cenários complexos. Ainda assim, a entrega atende ao objetivo do trabalho ao integrar reconhecimento de padrões, segmentação, morfologia, esqueletonização e grafos em um fluxo executável e documentado.

\begin{thebibliography}{9}

\bibitem{gonzalez}
R. C. Gonzalez e R. E. Woods.
\textit{Digital Image Processing}.
Pearson, 4ª edição, 2018.

\bibitem{szeliski}
R. Szeliski.
\textit{Computer Vision: Algorithms and Applications}.
Springer, 2ª edição, 2022.

\bibitem{opencv}
OpenCV Team.
\textit{OpenCV Documentation}.
Disponível em: \url{https://docs.opencv.org/}.

\bibitem{skimage}
scikit-image Developers.
\textit{scikit-image: Image processing in Python}.
Disponível em: \url{https://scikit-image.org/}.

\bibitem{networkx}
NetworkX Developers.
\textit{NetworkX Documentation}.
Disponível em: \url{https://networkx.org/}.

\end{thebibliography}

\end{document}
